% ======================================================================
%  WIT-0604151-EMDGL 解经本 — 共享导言区
%  编译：xelatex main.tex（两遍），工作目录 manuscript/
% ======================================================================

% ---------- 数学 ----------
\usepackage{amsmath,amssymb,amsthm,mathtools}
\usepackage{mathrsfs}
\usepackage{bm}

% ---------- 排版 ----------
\usepackage[margin=2.3cm,top=2.4cm,bottom=2.4cm]{geometry}
\usepackage{graphicx}
\usepackage{booktabs,array,longtable,multicol}
\usepackage{enumitem}
\setlist{itemsep=2pt,topsep=3pt,parsep=0pt}
\usepackage{xcolor}
\usepackage{xspace}
\usepackage{fancyhdr}
\usepackage[colorlinks=true,linkcolor=blue!55!black,urlcolor=teal!70!black,citecolor=blue!55!black]{hyperref}

% ---------- 图 ----------
\usepackage{tikz}
\usetikzlibrary{arrows.meta,positioning,calc,shapes.geometric,fit,backgrounds,decorations.pathreplacing}

% ---------- 颜色 ----------
\definecolor{kwblue}{RGB}{28,66,120}
\definecolor{kwgray}{RGB}{105,105,105}
\definecolor{kwred}{RGB}{150,40,40}
\definecolor{kwgreen}{RGB}{30,95,60}
\definecolor{kwboxbg}{RGB}{247,246,240}

% ---------- 标题样式 ----------
\usepackage{titlesec}
\titleformat{\section}{\Large\bfseries\color{kwblue}}{\thesection}{0.7em}{}
\titleformat{\subsection}{\large\bfseries\color{kwblue!85!black}}{\thesubsection}{0.6em}{}
\titlespacing*{\section}{0pt}{2.2ex plus .4ex}{1.2ex}
\titlespacing*{\subsection}{0pt}{1.8ex plus .3ex}{0.9ex}

% ---------- 页眉页脚 ----------
\pagestyle{fancy}
\fancyhf{}
\fancyhead[LE,RO]{\small\color{kwgray}\thepage}
\fancyhead[LO]{\small\color{kwgray}\nouppercase{\rightmark}}
\fancyhead[RE]{\small\color{kwgray}\nouppercase{\leftmark}}
\renewcommand{\headrulewidth}{0.3pt}
\setlength{\headheight}{14pt}
% 中文行内公式容易撑破行：给 TeX 一点弹性
\setlength{\emergencystretch}{2.2em}

% ---------- 定理类 ----------
\theoremstyle{definition}
\newtheorem{definition}{定义}[section]
\newtheorem{example}[definition]{例}
\newtheorem{remark}[definition]{评注}
\theoremstyle{plain}
\newtheorem{theorem}[definition]{定理}
\newtheorem{proposition}[definition]{命题}
\newtheorem{lemma}[definition]{引理}
\newtheorem{conjecture}[definition]{猜想}

% ---------- 解经专用宏 ----------
% 锚点：\anchor{src p.12} 或 \anchor{eq 3.7} 或 \anchor{§9.2}
\newcommand{\anchor}[1]{{\footnotesize\color{kwgray}[#1]}}
% provenance：A=原文明确说 B=原文暗示、我们重建 C=解经者判断
\newcommand{\prov}[1]{{\footnotesize\color{kwgray}\textsf{(#1)}}}
\newcommand{\pA}{\prov{A}\xspace}
\newcommand{\pB}{\prov{B}\xspace}
\newcommand{\pC}{\prov{C}\xspace}
% 术语首次出现：\term{英文}{中文}
\newcommand{\term}[2]{\textbf{#1}（#2）}
\newcommand{\kw}{\textsc{Kapustin--Witten}}
\newcommand{\Ldual}{{}^{L}G}
\newcommand{\LdualC}{{}^{L}G_{\mathbb{C}}}
\newcommand{\MH}{\mathcal{M}_H}
\newcommand{\MC}{\mathcal{M}(G,C)}
\newcommand{\gC}{\mathfrak{g}_{\mathbb{C}}}
\newcommand{\CC}{\mathbb{C}}
\newcommand{\PP}{\mathbb{P}}
\newcommand{\ZZ}{\mathbb{Z}}
\newcommand{\RR}{\mathbb{R}}
\newcommand{\HH}{\mathbb{H}}
\newcommand{\dd}{\mathrm{d}}
\newcommand{\ii}{\mathrm{i}}
\newcommand{\ee}{\mathrm{e}}
\newcommand{\Tr}{\operatorname{Tr}}
\newcommand{\Dmod}{\mathcal{D}}

% 骨架卡：\begin{kcard}{卡片标题} ... \end{kcard}
% （锚点写在卡片内第一行：\anchor{src p.12}）
\usepackage[most]{tcolorbox}
\newtcolorbox{kcard}[2][]{%
  enhanced, breakable,
  colback=white, colframe=kwblue!65, boxrule=0.6pt, arc=1.5pt,
  left=7pt, right=7pt, top=4pt, bottom=4pt,
  fonttitle=\bfseries\small, coltitle=white, colbacktitle=kwblue!65,
  attach boxed title to top left={xshift=6pt, yshift=-\tcboxedtitleheight/2},
  boxed title style={boxrule=0pt, arc=1pt},
  before skip=8pt, after skip=8pt,
  title={#2}, #1}

% 关键想法块：\begin{keyidea} ... \end{keyidea}
\newtcolorbox{keyidea}{%
  enhanced, breakable,
  colback=kwboxbg, colframe=kwboxbg, boxrule=0pt,
  borderline west={2.2pt}{0pt}{kwblue},
  left=9pt, right=7pt, top=5pt, bottom=5pt,
  fontupper=\small, before skip=8pt, after skip=8pt}

% 原文短引（≤200 字，锚点必附）：\begin{origquote} ... \end{origquote}
\newtcolorbox{origquote}{%
  enhanced, breakable,
  colback=white, colframe=white, boxrule=0pt,
  borderline west={1.2pt}{0pt}{kwgray},
  left=9pt, right=7pt, top=4pt, bottom=4pt,
  fontupper=\small\itshape, before skip=6pt, after skip=8pt}

% 警告块（环境版）：\begin{warn} ... \end{warn}
\newtcolorbox{warn}{%
  enhanced, breakable,
  colback=red!2!white, colframe=white, boxrule=0pt,
  borderline west={2.2pt}{0pt}{red!55!black},
  left=9pt, right=7pt, top=5pt, bottom=5pt,
  fontupper=\small, before skip=8pt, after skip=8pt}

% 「卡住点」标注
\newcommand{\stuck}[1]{{\color{kwred}\textbf{[卡住点]}\ #1}}
\newcommand{\caveat}[1]{{\color{kwred}\textbf{[注意]}\ #1}}

% 表头快捷
\newcommand{\tabhead}[1]{\textbf{\color{kwblue}#1}}

% ---------- 兜底：把圈号 ①–⑩ 与常用符号交给 CJK 字体 ----------
% 否则 xeCJK 会把它们留在 Latin 字体里，导致 "Missing character"（Fandol 收录 GB2312 圈号）。
\xeCJKDeclareCharClass{CJK}{"2460 -> "2473}
% 顺手把全角括号也强制到 CJK（个别表格里曾经漏判）
\xeCJKDeclareCharClass{CJK}{"FF08, "FF09}


% --- macros added after section drafting (were local to 21_layer1_geom.tex)
\newcommand{\Lw}{{}^{L}w}
\newcommand{\LR}{{}^{L}R}
